%Präambel
\documentclass[fontsize=12]{scrartcl}

\usepackage{cite}
\usepackage[english, ngerman]{babel}
\usepackage{mwe}
\usepackage{newfloat}
\usepackage[utf8]{inputenc}
\usepackage[autostyle=true,german=quotes]{csquotes}
\usepackage{lmodern,textcomp}
\usepackage[T1]{fontenc}
\usepackage{url}
\usepackage{amssymb}

\usepackage{amsmath}

\usepackage{tabulary}
\usepackage{tabularx}
\usepackage{subfig}
\usepackage{caption, booktabs}

\usepackage{float}

%der eigentliche Text
\begin{document}
	
	% Festlegung Art der Zitierung - Havardmethode: Abkuerzung Autor + Jahr
	\bibliographystyle{abbrvdin}
	
	%\maketitle
	
	\begin{center}
		
		\fontsize{14pt}{12pt}
		
		\textbf{Light Years Ahead}\\
		\textbf{- Gaming - }\\
		
	\end{center}
	
	\begin{verbatim}
		
		
		
		
	\end{verbatim}
	
	\begin{center}
		
		\fontsize{14pt}{12pt}
		
		\textbf{Design einer Software zum Erstellen von Sammelkartenspielen}
		
	\end{center}
	
	\begin{verbatim}
		
		
		
		
		
		
		
		
		
		
		
		
		
		
		
		
		
		
		
		
		
		
		
		
		
		
		
		
	\end{verbatim}
	
	\begin{flushright}
		
		Leipzig, Juli, 2025 \hfill vorgelegt von\\
		
		\bigskip
		
		Felix Cäsar Madorskiy\\
		Studiengang Informatik M.Sc.
		
	\end{flushright}
	
	\vfill
	
	\pagenumbering{gobble}
	
	\newpage
	
	%\newpage
	
	\setcounter{page}{2} 
	
	\newpage
	
	\section{Wer Ich Bin}
	
	Mein Name ist Felix und ich liebe Sammelkartenspiele -- egal ob Pokemon, Yu-Gi-Oh oder Magic: the Gathering. Doch bei all diesen Spielen bin ich immer wieder auf dasselbe Problem gestoßen: Die Herausgeber der Spiele drucken immer stärkere Karten, um ihr Produkt zu verkaufen, was ältere Karten unbrauchbar macht und geliebte Spielstrategieen unspielbar.\hfill\newline
	
	\noindent Ich will das ändern. Wenn du diese Vision teilst und Teil eines neuen Kapitels für Sammelkartenspiele sein möchtest -- lies weiter.
	
	\section{Projekt}
	
	Das Ziel ist das Kreieren einer Software, mit der man Sammelkartenspiele erstellen kann, ohne auch nur eine Zeile Code geschrieben zu haben. Was uns zusätzlich von der Konkurrenz unterscheidet, ist, dass ein integraler Bestandteil unserer Software die Pitchen-und-Bewerten Mechanik ist. Mit ihr können Spieler Karten, welche sie sich ausgedacht haben, der Spielergemeinschaft vorschlagen und die Spieler können dann gewichtet bewerten, ob die Karte in das Spiel aufgenommen wird. Gewichtet in dem Sinne, dass die Stimme von erfahrenen Spielern mehr wiegt. Das selbe System lässt Spieler auch vorschlagen, wie eine bereits vorhandene Karte verändert werden muss, um besser ins Spiele passen. Das verhindert, dass Karten in das Spiel aufgenommen werden, die für das Spiel desaströs sind und alte Karten und geliebte Spielstrategien obsolet machen. Das Ziel ist mithilfe der Software unser eigenes Sammelkartenspiel zu kreieren und dann später auch als physische, gedruckte Karten zu vermarkten aber wenn jemand anderes unsere Software benutzt, ist die Benutzung kostenlos aber wir kriegen 20\% von allen Verkäufen, welche (Sammel-)Kartenspiele mit unserer Software erzielen. Die integrierten Marketplaces zählen nicht dazu.\hfill\newline
	
	\noindent Ferner soll die Software drei integrierte Marketplaces haben, einen für Bilder für die Kartengestaltung, einen für Musik, die beim Ausspielen oder bei der Aktivierung von Karteneffekten abgespielt wird und einen für Merchandise rund um die Kartenspiele und darüber hinaus. Wir würden 20\% des Verkaufspreises als Gebühr bekommen.\hfill\newline
	
	\noindent Schließlich soll die Software es ermöglichen verschiedene Spielmodi zu erstellen und direkt in der Software Wettbewerbe zu veranstalten, sowohl kostenlos als auch mit monetären Prämien für die Sieger, sodass sowohl professionelle Spieler als auch professionelle Streamer in unserer Software Karrierechancen sehen. Auch hier würden wir 20\% Provision haben.
	
	\section{Team}
	
	Wir sind bereits zu fünft. Außer mir, dem Visionär hinter dem Vorhaben, sind in unserem Team drei Programmierer und ein Musiker. Von den Programmierern macht einer Frontend, einer Backend und einer bringt sich allgemein ein. Der Programmierer, welcher für das Frontend zuständig ist, macht auch Design. Der Musiker macht Ingame Sounds und Spielmusik. Leider bringen sich alle nur 4 Stunden pro Woche ein, was sie für das EXIST-Gründerstipendium ausschließt, für welches man am Projekt in Vollzeit arbeiten muss.
	
	\section{Was Ich Suche}
	
	Ich suche eine*n Informatiker*in als Co-Gründer*in für die Bewerben auf das EXIST-Gründerstipendium. Das Exist-Gründerstipendium unterstützt Gründer*innen dabei, ihr Projekt für ein Jahr zu entwickeln, finanziert mit 2.500€ pro Person/Monat (zzgl. Zuschläge), vorausgesetzt der/die Gründer*in hat mindestens einen Bachelorabschluss. Das EXIST-Gründerstipendium ist ein Vollstipendium und muss nicht zurückgezahlt werden.\hfill\newline
	
	\noindent Deine Aufgaben werden sein:
	
	\begin{itemize}
		
		\item Programmieren eines rudimentären Launchers
		
		\item Programmieren rudimentärer Marketplaces
		
		\item Programmieren des Feld-Editors

		\item Programmieren des Karten-Editors
		
		\item Programmieren eines rudimentären Supports für das Einbinden von Sounds/Musik/Effekten, die abgespielt werden, wenn ein Spieler auf das Feld/die Karte drückt oder die Karte ausspielt bzw. deren Effekt benutzt, auch abhängig davon, von wo die Karte ausgespielt wird bzw. wo/wann ihr Effekt benutzt wird, z.B. aus der Hand oder dem Ablagestapel
		
	\end{itemize}
	
	\noindent Unser Tech-Stack ist:
	
	\begin{itemize}
		
		\item Discord
		
		\item Git
		
		\item Godot (die Game Engine)
		
		\item GDScript
		
		\item AWS/Rust
		
		\item C\#
		
	\end{itemize}
	
	\noindent Ich freue mich auf deine Antwort.
	
	\section{Kontakt}
	
	E-Mail: felix.caesar.madorskiy@gmail.com\\
	LinkedIn: Felix Cäsar Madorskiy\\
	Discord: Tsesar\#0749 bzw. \_tsesar oder FXJack
	
	% Literaturliste soll im Inhaltsverzeichnis auftauchen
	%\addcontentsline{toc}{section}{Literaturverzeichnis}
	
	% Literaturliste endgueltig anzeigen
	%\bibliography{Smile_Antrag.bib}
	
\end{document}